%-------------------------------------------------------------------------------
\section{Background and Motivation}
\label{sec:background}

\subsection{Keyspace Partitioning in Distributed In-Memory Stores}
\label{sec:partitioning}

A foundational design choice in any distributed data store is the mechanism by which the keyspace is divided among nodes~\cite{slicer,shard-manager}. Partitioning determines not only where a given record resides at any point in time, but also how the system responds to membership changes: which records must move, how they are located, and how quickly they can be transferred~\cite{akkio}. Across systems, two broad families of partitioning scheme are in use---static hash-slot assignment and consistent hashing on a token ring~\cite{karger-consistent-hashing,dynamodb}---and both give rise to the same underlying memory-layout problem that SAMM addresses.

\paragraph{Hash-slot partitioning.}
In this approach, the keyspace is divided into a fixed number of discrete slots, each assigned to exactly one node. Keys are mapped to slots via a hash function applied to the key. Redis Cluster, which uses 16,384 slots assigned via CRC16~\cite{redisclusterconfig}, is the canonical example of this model.

\paragraph{Token-ring partitioning.}
Alternatively, nodes are assigned positions on a circular hash ring and are responsible for all keys whose hash falls within their token range. Apache Cassandra~\cite{cassandra-paper}, which supports multiple virtual token ranges per node to improve load balance, and Amazon Dynamo~\cite{dynamodb} and DynamoDB~\cite{dynamodb-atc22}, which split and reassign hash ranges automatically as tables grow, follow this approach.

\paragraph{Data layout does not follow the logical data organization.}
In both schemes, rebalancing requires gathering all records belonging to a partition unit---a slot or a token range---from wherever the storage layer has placed them and streaming them to a new owner~\cite{rocksteady,netmigrate}. For disk-based systems such as Cassandra, sorted SSTable files~\cite{bigtable} provide natural co-location within a token range, reducing streaming to sequential reads. For in-memory systems such as Redis, no such structure exists~\cite{mica,memc3}. A general-purpose heap allocator interleaves records from all slots without regard to slot membership, so gathering a single slot's records requires a full hash-index traversal followed by random pointer-chased reads across the heap. It is precisely this regime---in-memory, hash-indexed, mutation-heavy~\cite{atikoglu-kv-workload,yang-twitter-cache,scaling-memcache}---that SAMM is designed to address.

\subsection{The Memory-Layout Bottleneck}
\label{sec:layout-bottleneck}

General-purpose allocators such as jemalloc~\cite{jemalloc} and TCMalloc~\cite{tcmalloc} optimise for throughput and low fragmentation across arbitrary allocation patterns~\cite{rumble-log-memory}. They organise free memory into size-class bins and per-thread arenas~\cite{hoard_paper}. The physical placement of two allocations is determined by availability within those bins, not by any semantic relationship between the objects. As a result, records belonging to the same slot are interleaved with records from all other slots throughout the heap.

During slot migration, the server must traverse its global hash table, identify all keys in the target slot, and dereference each value pointer to retrieve the associated record~\cite{rocksteady}. These pointer dereferences generate random accesses across the heap, thrashing CPU caches and precluding vectorised or DMA-assisted bulk copies~\cite{kvdirect,lite-rdma}. The migration throughput is therefore bounded not by network bandwidth but by memory-access latency---a fundamental impedance mismatch between the application's logical model and the allocator's physical layout.

\subsection{Motivating Example}
\label{sec:motivation-experiment}

\textcolor{red}{[OPEN: ``preview'' results]}
